
Now, by VI 4.8, @@M0@@ is the @@M1@@-algebra generated by the @@M2@@, subject to the relations deduced from:
@@M3@@
Thus @@M4@@, where @@M5@@ is the ideal generated by the @@M6@@ for @@M7@@, hence the @@M8@@-ideal generated by @@M9@@. Consequently, @@M10@@. The proof is completed by the lemma:

\sgaStableTarget{sga6.e09.lem.4-2}\textbf{Lemma 4.2.} \uline{Let} @@M11@@ \uline{be a noetherian scheme}, @@M12@@ \uline{a parabolic subgroup of a reductive} @@M13@@\uline{-group}, \uline{and} @@M14@@ \uline{a torsor under the unipotent radical} @@M15@@ \uline{of} @@M16@@; \uline{then one has an isomorphism}:
@@M17@@
Indeed, there is a sequence of subgroup schemes of @@M18@@:
@@M19@@
closed in @@M20@@, finite because @@M21@@ is noetherian, and such that, for every @@M22@@, @@M23@@ is isomorphic to an @@M24@@-vector group (SGA 3 XXVI 2.1). Consequently, @@M25@@ factors through a sequence of torsors under vector groups, whence the result by 3.4.

\sgaStableTarget{sga6.e09.cor.4-3}\textbf{Corollary 4.3.} \uline{Let} @@M26@@ \uline{be a noetherian scheme}, @@M27@@. \uline{Then the canonical homomorphism} @@M28@@ \uline{is an isomorphism}.

This follows immediately from 4.1 applied to @@M29@@, since then @@M30@@.

